% Options for packages loaded elsewhere
\PassOptionsToPackage{unicode}{hyperref}
\PassOptionsToPackage{hyphens}{url}
\PassOptionsToPackage{dvipsnames,svgnames,x11names}{xcolor}
%
\documentclass[
  11pt,
]{article}
\usepackage{amsmath,amssymb}
\usepackage{iftex}
\ifPDFTeX
  \usepackage[T1]{fontenc}
  \usepackage[utf8]{inputenc}
  \usepackage{textcomp} % provide euro and other symbols
\else % if luatex or xetex
  \usepackage{unicode-math} % this also loads fontspec
  \defaultfontfeatures{Scale=MatchLowercase}
  \defaultfontfeatures[\rmfamily]{Ligatures=TeX,Scale=1}
\fi
\usepackage{lmodern}
\ifPDFTeX\else
  % xetex/luatex font selection
\fi
% Use upquote if available, for straight quotes in verbatim environments
\IfFileExists{upquote.sty}{\usepackage{upquote}}{}
\IfFileExists{microtype.sty}{% use microtype if available
  \usepackage[]{microtype}
  \UseMicrotypeSet[protrusion]{basicmath} % disable protrusion for tt fonts
}{}
\makeatletter
\@ifundefined{KOMAClassName}{% if non-KOMA class
  \IfFileExists{parskip.sty}{%
    \usepackage{parskip}
  }{% else
    \setlength{\parindent}{0pt}
    \setlength{\parskip}{6pt plus 2pt minus 1pt}}
}{% if KOMA class
  \KOMAoptions{parskip=half}}
\makeatother
\usepackage{xcolor}
\usepackage[margin=1in]{geometry}
\setlength{\emergencystretch}{3em} % prevent overfull lines
\providecommand{\tightlist}{%
  \setlength{\itemsep}{0pt}\setlength{\parskip}{0pt}}
\setcounter{secnumdepth}{-\maxdimen} % remove section numbering
\DeclareUnicodeCharacter{220E}{\ensuremath{\square}}
\ifLuaTeX
  \usepackage{selnolig}  % disable illegal ligatures
\fi
\IfFileExists{bookmark.sty}{\usepackage{bookmark}}{\usepackage{hyperref}}
\IfFileExists{xurl.sty}{\usepackage{xurl}}{} % add URL line breaks if available
\urlstyle{same}
\hypersetup{
  pdfauthor={Author information to be supplied before submission},
  colorlinks=true,
  linkcolor={Maroon},
  filecolor={Maroon},
  citecolor={Blue},
  urlcolor={Blue},
  pdfcreator={LaTeX via pandoc}}

\title{Alphabet-Davenport Normal Forms for Multi-Tag Quantum Compilation}
\author{Author information to be supplied before submission}
\date{Review source - 25 August 2026}

\begin{document}
\maketitle

\hypertarget{abstract}{%
\section{Abstract}\label{abstract}}

Finite-support normal forms make exact quantum-compiler optimization
enumerable, but an ambient-rank bound can hide the finer combinatorial
object that controls deletion. Let \(H\) be a finite abelian signature
group and \(A\subseteq H\) an allowed alphabet. Write
\(\operatorname{zsf}(H; A)\) for the largest length of a word over \(A\)
with no nonempty zero-sum subword. We prove that every instance of a
deletion-closed compiler grammar has an exact optimum of support at most
\(\operatorname{zsf}(H; A)\), provided each constrained generator has
nonzero total signature and every zero-signature deletion preserves
feasibility without increasing cost. For \(H=\mathbb F_2^d\), this
implies \(\operatorname{zsf}(H; A)\le d\), with equality whenever \(A\)
contains a basis.

We instantiate the theorem in an arbitrary-block MultiTag-TARE grammar.
With \(b\ge2\) ordered blocks, \(s\ge0\) shared Tags, minimum frame
coefficient \(\mu\), and Restore coefficient \(t_R\), changing one local
letter increases the \(b\)-way Restore functional by at most \(b-1\),
and this bound is sharp. Consequently, throughout the sufficient cone

\[
\mu\ge (b-1)t_R,
\]

every instance has an optimum satisfying

\[
|\operatorname{supp}(R)|
\le \operatorname{zsf}(H_R; A_R)
\le \operatorname{rank}(H_R)
\le s+1
\]

for every frame \(R\). The one-Tag, three-block R6M specialization is
sharply intrinsic, with \(\kappa_{\mathrm{R6M}}=2\). We make no
necessity claim for multiple Tags or outside the sufficient cone. The
result separates the alphabet-sensitive certificate, the compiler
semantics that make deletion sound, and the objective inequality that
makes deletion profitable.

\textbf{Keywords:} quantum compilation; block encoding; exact normal
forms; zero-sum sequences; sparse optimization

\hypertarget{introduction}{%
\section{1. Introduction}\label{introduction}}

TARE is an upstream block-encoding construction that exposes frame, Tag,
branch, and Restore choices. The present work studies exact normal forms
for those auxiliary compiler choices; it does not claim the underlying
block-encoding primitive {[}1{]}.

A standard support proof attaches a finite binary syndrome to every
active coordinate. When support exceeds the syndrome dimension, linear
dependence produces a removable zero-syndrome subset. The dimension
obtained in this way can look like an intrinsic property of the
compiler, although the proof uses only two ingredients: a sufficiently
long signature word contains a zero-sum subword, and deleting that
subword is non-increasing. The primitive threshold is therefore
alphabet-sensitive rather than rank-sensitive.

This distinction matters for three reasons. First, it extends the
deletion argument from binary vector spaces to finite abelian signature
groups. Second, even in \(\mathbb F_2^d\), the alphabet realized by a
compiler can have a smaller zero-sum-free ceiling than the ambient
dimension. Third, it keeps the combinatorial certificate separate from
the objective inequality that licenses the corresponding edit.

Our contributions are:

\begin{enumerate}
\def\labelenumi{\arabic{enumi}.}
\tightlist
\item
  an alphabet-Davenport deletion theorem for deletion-closed
  optimization grammars with persistent soundness and dominance
  assumptions;
\item
  the binary rank corollary, including exactness for the deletion
  language when a basis word is allowed;
\item
  an arbitrary-block MultiTag-TARE instantiation with exact local
  Restore sensitivity \(b-1\) and sufficient cone \(\mu\ge(b-1)t_R\);
\item
  a sharp frozen R6M control with intrinsic support two; and
\item
  an explicit boundary between certificate support, intrinsic support,
  and physical resource claims.
\end{enumerate}

\hypertarget{an-alphabet-restricted-zero-sum-invariant}{%
\section{2. An alphabet-restricted zero-sum
invariant}\label{an-alphabet-restricted-zero-sum-invariant}}

Let \(H\) be a finite abelian group written additively, and let
\(A\subseteq H\) be finite. A word over \(A\) is \emph{zero-sum-free} if
none of its nonempty subwords has sum zero. Define

\[
\operatorname{zsf}(H; A)
=\max\{|W|:W\text{ is a zero-sum-free word over }A\}.
\]

Words permit repetition, matching the multiplicity of compiler
coordinates. For \(A=H\setminus\{0\}\), the invariant equals \(D(H)-1\),
where \(D(H)\) is the classical Davenport constant. Smaller alphabets
can yield strictly smaller values. We use explicit notation to avoid
conflating this quantity with other weighted or restricted zero-sum
conventions.

\hypertarget{the-deletion-theorem}{%
\section{3. The deletion theorem}\label{the-deletion-theorem}}

Fix an optimization grammar. For each constrained generator \(R\), every
active coordinate \(q\) carries a signature \(v_q\in A_R\subseteq H_R\).
Assume the following conditions hold initially and after every admitted
deletion:

\begin{enumerate}
\def\labelenumi{\arabic{enumi}.}
\tightlist
\item
  \textbf{Nonzero total.} Every feasible constrained generator satisfies
  \(\sum_q v_q\ne0\).
\item
  \textbf{Deletion closure and soundness.} Replacing \(R\) by the
  identity on any nonempty zero-sum subword remains inside the admitted
  grammar and preserves every semantic constraint represented by the
  signatures.
\item
  \textbf{Persistent deletion dominance.} The same operation does not
  increase the objective.
\end{enumerate}

\textbf{Theorem 1 (alphabet-Davenport normal form).} Every admitted
instance has an exact optimum for which

\[
|\operatorname{supp}(R)|\le \operatorname{zsf}(H_R; A_R)
\]

for every constrained generator \(R\).

\textbf{Proof.} Begin with an optimum. If the signature word of \(R\) is
longer than \(\operatorname{zsf}(H_R; A_R)\), it contains a nonempty
zero-sum subword. The nonzero total prevents that subword from being the
whole word. Delete its coordinates. Closure and soundness preserve
feasibility, while dominance preserves optimality. Support strictly
decreases. Because all three assumptions persist, the process terminates
with support at most \(\operatorname{zsf}(H_R; A_R)\). Repeating the
argument for each constrained generator gives a simultaneous normal
form. ∎

The abstract statement localizes the compiler-specific burden: one must
define the signature map, prove nonzero total and deletion closure, and
account for the objective change. The zero-sum threshold alone supplies
none of those facts.

\hypertarget{binary-rank-as-a-corollary}{%
\section{4. Binary rank as a
corollary}\label{binary-rank-as-a-corollary}}

Let \(H=\mathbb F_2^d\). Every word longer than \(d\) is linearly
dependent and therefore contains a nonempty zero-XOR subword. Hence

\[
\operatorname{zsf}(\mathbb F_2^d; A)\le d
\]

for every \(A\subseteq\mathbb F_2^d\). If \(A\) contains a basis, the
basis word is zero-sum-free, so equality holds.

\textbf{Corollary 2.} A rank-\(d\) deletion bound in an elementary
binary signature system is exact for the stated deletion language when a
basis word is allowed. It becomes an intrinsic compiler lower bound only
when an independent compiler witness proves that support below \(d\) is
impossible.

\hypertarget{multitag-tare-grammar}{%
\section{5. MultiTag-TARE grammar}\label{multitag-tare-grammar}}

Fix \(b\ge2\) ordered two-target blocks and \(s\ge0\) shared Tag Paulis
\(S_1, \ldots, S_s\). Each frame Pauli \(R\) has an anticommuting
partner \(R'\). At every active coordinate, define

\[
v_q=\bigl(\langle R_q, R'_q\rangle,
\langle S_{1q}, R_q\rangle, \ldots,
\langle S_{sq}, R_q\rangle\bigr)
\in\mathbb F_2^{s+1},
\]

where \(\langle\cdot, \cdot\rangle\) is the local symplectic commutation
form. The XOR of the first components is one, so the total signature is
nonzero.

The structural objective charges every active frame coordinate by a
coefficient of at least \(\mu\), permits arbitrary nonnegative
Tag-support charges, and adds \(t_R\ge0\) times the local Restore
functional

\[
F_b(a_1, \ldots, a_b)=
\begin{cases}
1, & a_1=\cdots=a_b\ne I, \\
|\{i:a_i\ne I\}|, & \text{otherwise}.
\end{cases}
\]

\hypertarget{exact-local-restore-sensitivity}{%
\section{6. Exact local Restore
sensitivity}\label{exact-local-restore-sensitivity}}

\textbf{Lemma 3.} Replacing one argument of \(F_b\) increases its value
by at most \(b-1\), and the bound is attained.

\textbf{Proof.} Away from the all-equal nonidentity state, \(F_b\) is
the ordinary nonidentity Hamming weight, so one replacement raises it by
at most one. Entering the exceptional state lowers the value to one.
Leaving that state while keeping the changed letter nonidentity changes
the value from one to \(b\), giving the exact increase \(b-1\). ∎

\hypertarget{multitag-normal-form}{%
\section{7. MultiTag normal form}\label{multitag-normal-form}}

Let \(A_R\) be the alphabet actually realized by the partner/Tag
signatures of frame \(R\), and let \(H_R=\langle A_R\rangle\).

\textbf{Theorem 4.} If \(\mu\ge(b-1)t_R\), every admitted MultiTag-TARE
instance has an exact optimum satisfying

\[
|\operatorname{supp}(R)|
\le \operatorname{zsf}(H_R; A_R)
\le \operatorname{rank}(H_R)
\le s+1
\]

for every frame \(R\).

\textbf{Proof.} A zero-signature subword preserves partner
anticommutation and every Tag label; the Tags remain fixed. Deleting
each selected coordinate refunds at least \(\mu\) in frame cost and can
add at most \((b-1)t_R\) in Restore cost by Lemma 3. The deletion is
therefore non-increasing in the stated cone. Its effect stays within the
same grammar, and the signature invariants persist, so Theorem 1 applies
iteratively. Finally, \(H_R\) is an elementary binary subgroup of
\(\mathbb F_2^{s+1}\), which gives the rank inequalities. ∎

The alphabet ceiling can be strictly smaller than realized rank when the
allowed alphabet contains no basis word of its generated subgroup. This
is the instance-adaptive refinement that an ambient dimension alone
cannot express.

\hypertarget{sharp-r6m-specialization}{%
\section{8. Sharp R6M specialization}\label{sharp-r6m-specialization}}

For the frozen one-Tag, three-block R6M grammar,

\[
b=3, \qquad s=1, \qquad \mu=2, \qquad t_R=1.
\]

The objective lies on the sufficient-cone boundary. The signature
alphabet realizes a basis of \(\mathbb F_2^2\), so the deletion
certificate ceiling is two. An independent all-size normalization proves
support at most two, and an exact support-one obstruction proves
necessity. Thus

\[
\kappa_{\mathrm{R6M}}=2.
\]

Only this frozen specialization receives an intrinsic sharpness claim.

\hypertarget{relation-to-prior-work}{%
\section{9. Relation to prior work}\label{relation-to-prior-work}}

Sparse optimal solutions {[}2{]}, Davenport constants, restricted
zero-sum invariants {[}3{]}, finite-field dependence, Pauli symplectic
algebra, and exact synthesis are established areas. In particular,
general sparse integer optimization provides objective-independent
support bounds, while zero-sum theory provides sequence thresholds for
selected alphabets and weights. Those results own the generic ideas of
sparse optima and zero-sum thresholds.

The residual contribution here is their compiler-specific conjunction: a
TARE-derived semantic signature, exact arbitrary-block Restore
sensitivity, a sufficient objective cone, an alphabet-sensitive support
normal form, and a separately sharp R6M control.

\hypertarget{reproducibility-and-limitations}{%
\section{10. Reproducibility and
limitations}\label{reproducibility-and-limitations}}

Finite-group controls, binary basis obstructions, and the exact local
Restore sensitivity have been checked independently for bounded
instances. The proofs above carry the all-size authority.

The theorem applies only when zero-sum deletion is grammar-closed,
semantics-preserving, and non-increasing after every prior deletion.
Computing \(\operatorname{zsf}(H; A)\) may itself be difficult. The
MultiTag result applies only to the explicit grammar and objective.
Outside \(\mu\ge(b-1)t_R\), the proof is silent; it does not imply that
larger support is necessary. General multi-Tag sharpness remains open.
Structural support is not physical T count, runtime, circuit depth,
qubit count, or quantum advantage.

\hypertarget{data-and-code-availability}{%
\section{Data and code availability}\label{data-and-code-availability}}

The exact bounded controls and deterministic verification scripts
supporting the finite-group, basis-obstruction, and Restore-sensitivity
checks accompany the submission source. They are corroborative: the
displayed proofs carry the all-size claims. A permanent archival
identifier must be inserted before the final journal upload.

\hypertarget{conclusion}{%
\section{11. Conclusion}\label{conclusion}}

Ambient rank is not the primitive deletion threshold. The exact
threshold is the longest zero-sum-free word expressible by the realized
signature alphabet. Rank remains a convenient binary upper bound and is
exact for the deletion language when a basis is realized. Only a
separate compiler witness can turn that certificate ceiling into
intrinsic support.

For MultiTag-TARE, block count controls deletion profitability through
the \(b-1\) Restore penalty, while the realized alphabet controls
sufficient support. The frozen R6M case aligns certificate and intrinsic
complexity at two. Beyond that case, the result is a strong normal form,
not a necessity theorem.

\hypertarget{references}{%
\section{References}\label{references}}

\begin{enumerate}
\def\labelenumi{\arabic{enumi}.}
\tightlist
\item
  N. Schillo, A. Sturm, and R. Quay, ``Block Encoding Linear
  Combinations of Pauli Strings Using the Stabilizer Formalism,''
  arXiv:2601.05740 (2026).
\item
  I. Aliev, J. A. De Loera, F. Eisenbrand, T. Oertel, and R. Weismantel,
  ``The Support of Integer Optimal Solutions,'' \emph{SIAM Journal on
  Optimization} \textbf{28}, 2152-2157 (2018). DOI: 10.1137/17M1162792
\item
  G. Wang, ``The universal zero-sum invariant and weighted zero-sum for
  infinite abelian groups,'' \emph{Communications in Algebra}
  \textbf{53}(4), 1581-1599 (2025). DOI: 10.1080/00927872.2024.2418017
\end{enumerate}

\end{document}
